\documentclass[conference]{IEEEtran}
\usepackage[utf8]{inputenc}
\usepackage[spanish,es-tabla]{babel}
\usepackage{amsmath,amssymb,amsfonts}
\usepackage{algorithmic}
\usepackage{graphicx}
\usepackage{textcomp}
\usepackage{xcolor}
\usepackage{booktabs}
\usepackage{cite}
\usepackage{url}
\usepackage{hyperref}

\hypersetup{
    colorlinks=true,
    linkcolor=blue,
    citecolor=blue,
    urlcolor=blue
}

\begin{document}

\title{Sistema H\'{\i}brido Neuro-Simb\'olico de Visi\'on por Computador y Programaci\'on por Restricciones para la Resoluci\'on de KenKen}

\author{\IEEEauthorblockN{Proyecto de Investigaci\'on Aplicada}
\IEEEauthorblockA{\textit{T\'opicos en Ciencias de la Computaci\'on} \\
\textit{Universidad Peruana de Ciencias Aplicadas (UPC)}\\
Lima, Per\'u}
}

\maketitle

\begin{abstract}
Este art\'{\i}culo presenta una arquitectura integral neuro-simb\'olica para la resoluci\'on automatizada de acertijos aritm\'etico-l\'ogicos KenKen a partir de im\'agenes sin procesar. El sistema articula una fase sensorial de visi\'on computacional (correcci\'on de perspectiva, segmentaci\'on de cuadr\'{\i}cula y reconocimiento \'optico de caracteres con probabilidades de confusi\'on) con un motor de Programaci\'on por Restricciones (CP) sustentado en CP-SAT de Google OR-Tools. Formulamos el problema formalmente como un CSP $\langle X, D, C \rangle$, implementando y contrastando emp\'{\i}ricamente: (1) modelado aritm\'etico intensional con descomposici\'on de productos y divisi\'on reificada, (2) restricciones globales extensionales de tabla con garant\'{\i}a de Consistencia de Arco Generalizada (GAC), y (3) inferencia conjunta neuro-simb\'olica (MAP) para recuperaci\'on robusta ante fallos de lectura OCR. En evaluaciones sistem\'aticas sobre \'ordenes $n \in [3, 9]$, el solver resuelve tableros de hasta $9 \times 9$ ($1.96 \times 10^{77}$ estados brutos) en menos de 21~ms en el nodo ra\'{\i}z con cero ramas de backtracking, demostrando una reducci\'on del 39.5\% en tiempo de resoluci\'on mediante restricciones de tabla en $n=9$.
\end{abstract}

\begin{IEEEkeywords}
Programaci\'on por Restricciones, KenKen, Inferencia Neuro-Simb\'olica, Restricciones de Tabla, GAC, CP-SAT, Visi\'on por Computador.
\end{IEEEkeywords}

\section{Introducci\'on}
El acertijo KenKen, inventado en 2004 por Tetsuya Miyamoto, generaliza el concepto de Cuadrado Latino al imponer restricciones aritm\'eticas sobre subconjuntos conexos de celdas denominados \textit{jaulas} (\textit{cages}). A diferencia de Sudoku, la topolog\'{\i}a irregular de las jaulas y la multiplicidad de operadores ($=, +, -, \times, \div$) provocan una explosi\'on combinatoria exponencial en el espacio de estados no restringido $\mathcal{O}(n^{n^2})$.

El desaf\'{\i}o t\'ecnico se intensifica cuando el tablero debe inferirse directamente desde im\'agenes f\'{\i}sicas o capturas de c\'amara: peque\~nas ambig\"uedades o ruidos en el reconocimiento \'optico de caracteres (OCR) tornan el CSP r\'{\i}gido en insatisfactible si se asumen lecturas deterministas ingenuas. Para solventar este paradigma, proponemos una formulaci\'on neuro-simb\'olica h\'{\i}brida donde la visi\'on provee distribuciones de probabilidad sobre hip\'otesis visuales y el solver de CP resuelve de forma conjunta la consistencia l\'ogica y la optimizaci\'on sensorial v\'{\i}a reificaci\'on condicional.

\section{Arquitectura del Pipeline End-to-End}
El flujo de procesamiento se estructura en tres etapas secuenciales con lazo de retroalimentaci\'on correctivo:
\begin{enumerate}
    \item \textbf{Preprocesamiento y Rectificaci\'on Geom\'etrica:} Detecci\'on del contorno cuadrangular dominante de la grilla mediante transformadas morfol\'ogicas y rectificaci\'on de perspectiva mediante transformaci\'on proyectiva homogr\'afica.
    \item \textbf{Detecci\'on Topol\'ogica y OCR:} Segmentaci\'on de bordes de jaula, extracci\'on de las regiones de texto/glifos y clasificaci\'on con matrices de confusi\'on para obtener un ranking de candidatos $(T_{c,k}, \text{op}_{c,k}, P_{c,k})$.
    \item \textbf{Motor de Constraint Programming y Fallback:} Resoluci\'on inicial r\'apida con hip\'otesis \textit{top-1}. En caso de insatisfactibilidad (\texttt{INFEASIBLE}), el sistema activa de forma totalmente autom\'atica el modelo de inferencia conjunta neuro-simb\'olica (Variante C).
\end{enumerate}

% Inclusión del modelado formal de CP, complejidad y resultados experimentales
% ==============================================================================
% SECCIÓN DE CONSTRAINT PROGRAMMING (FORMATO IEEEtran)
% ==============================================================================

\section{Modelado Formal de Constraint Programming}
\label{sec:cp_modeling}

El acertijo l\'ogico KenKen de orden $n$ se formaliza rigurosamente como un Problema de Satisfacci\'on de Restricciones (CSP) definido por la tupla can\'onica $\mathcal{P} = \langle X, D, C \rangle$, donde $X$ es el conjunto de variables de decisi\'on, $D$ representa sus dominios asociados y $C$ es el conjunto de restricciones globales, aritm\'eticas y reificadas.

\subsection{Variables de Decisi\'on y Dominios}
Para una cuadr\'{\i}cula bidimensional de dimensiones $n \times n$, se define la matriz de variables de decisi\'on discretas:
\begin{equation}
    X = \{x_{i,j} \mid 1 \le i \le n, \; 1 \le j \le n\}
\end{equation}
donde cada variable entera $x_{i,j}$ denota el d\'{\i}gito asignado a la celda ubicada en la fila $i$ y columna $j$. El dominio uniforme para cada celda corresponde al conjunto finito de enteros:
\begin{equation}
    D(x_{i,j}) = \{1, 2, \dots, n\}, \quad \forall x_{i,j} \in X
\end{equation}
El problema contiene exactamente $n^2$ variables principales. La colecci\'on de jaulas $\mathcal{C} = \{c_1, \dots, c_m\}$ conforma una partici\'on exacta del tablero:
\begin{equation}
    \bigcup_{c \in \mathcal{C}} V_c = X \quad \land \quad V_a \cap V_b = \emptyset, \quad \forall a \ne b
\end{equation}
donde cada jaula $c = \langle V_c, T_c, \text{op}_c \rangle$ est\'a definida por un subconjunto conexo de celdas $V_c \subseteq X$, un objetivo num\'erico entero $T_c \in \mathbb{N}$ y un operador aritm\'etico $\text{op}_c \in \{=, +, -, \times, \div\}$.

\subsection{Restricciones Globales de Tablero (Cuadrado Latino)}
Para satisfacer la propiedad fundamental de no repetici\'on de valores en l\'{\i}neas ortogonales, se imponen restricciones globales de desigualdad sobre cada fila y cada columna mediante el algoritmo de filtrado de consistencia hiperarco de R\'egin~\cite{regin1994filtering}:
\begin{align}
    \text{AllDifferent}(\{x_{i,1}, x_{i,2}, \dots, x_{i,n}\}), & \quad \forall i \in \{1, \dots, n\} \label{eq:alldiff_row} \\
    \text{AllDifferent}(\{x_{1,j}, x_{2,j}, \dots, x_{n,j}\}), & \quad \forall j \in \{1, \dots, n\} \label{eq:alldiff_col}
\end{align}

Adicionalmente, se modelan restricciones globales redundantes (implicadas) basadas en la conservaci\'on de la suma triangular de Gauss:
\begin{equation}
    \sum_{j=1}^n x_{i,j} = \frac{n(n+1)}{2}, \quad \sum_{i=1}^n x_{i,j} = \frac{n(n+1)}{2}
    \label{eq:redundant_sum}
\end{equation}
Si bien (\ref{eq:redundant_sum}) es l\'ogicamente subsumida por (\ref{eq:alldiff_row})--(\ref{eq:alldiff_col}), su inclusi\'on aporta propagaci\'on anticipada de cotas lineales (\textit{bounds propagation}) en tableros de gran escala.

\subsection{Modelado de Jaulas: Variante A (Aritm\'etica Intensional)}
La formulaci\'on intensional descompone las relaciones matem\'aticas de cada jaula $c$ seg\'un su operador can\'onico:

\begin{enumerate}
    \item \textbf{Identidad ($=$):} Para jaulas unarias ($|V_c| = 1$), fijaci\'on directa de dominio:
    \begin{equation}
        v_1 = T_c, \quad v_1 \in V_c
    \end{equation}

    \item \textbf{Suma ($+$):} Restricci\'on af\'{\i}n lineal global sobre el vector de variables:
    \begin{equation}
        \sum_{v \in V_c} v = T_c
    \end{equation}

    \item \textbf{Multiplicaci\'on ($\times$):} Para jaulas asociativas con $|V_c| \ge 2$, se realiza una descomposici\'on encadenada mediante variables enteras auxiliares de producto $p_k \in [1, n^{k+1}]$ implementada v\'{\i}a \texttt{AddMultiplicationEquality}:
    \begin{align}
        p_1 &= v_1 \cdot v_2 \\
        p_k &= p_{k-1} \cdot v_{k+1}, \quad \forall k \in \{2, \dots, |V_c|-2\} \\
        T_c &= p_{|V_c|-2} \cdot v_{|V_c|}
    \end{align}

    \item \textbf{Diferencia Absoluta ($-$):} Restricci\'on binaria no lineal ($|V_c| = 2$):
    \begin{equation}
        |v_1 - v_2| = T_c \iff \text{AddAbsEquality}(T_c, v_1 - v_2)
    \end{equation}

    \item \textbf{Divisi\'on Reificada ($\div$):} Dado que en la imagen el orden relativo de los operandos no est\'a predefinido, se modela mediante \textbf{reificaci\'on l\'ogica condicional} introduciendo una variable booleana de direcci\'on $b_{\text{dir}} \in \{0, 1\}$ con la primitiva \texttt{OnlyEnforceIf}:
    \begin{align}
        b_{\text{dir}} &\implies v_1 = T_c \cdot v_2 \label{eq:div_reif_1} \\
        \neg b_{\text{dir}} &\implies v_2 = T_c \cdot v_1 \label{eq:div_reif_2}
    \end{align}
\end{enumerate}

\subsection{Modelado de Jaulas: Variante B (Restricciones Globales de Tabla)}
En la Variante B, cada jaula se formula de manera extensional mediante la enumeraci\'on previa del cat\'alogo de asignaciones v\'alidas $\mathcal{T}_c \subset \{1, \dots, n\}^{|V_c|}$, aplicando \texttt{AddAllowedAssignments}:
\begin{equation}
    \text{AllowedAssignments}(V_c, \mathcal{T}_c)
    \label{eq:table_constraint}
\end{equation}
Durante la construcci\'on previa de $\mathcal{T}_c$, el algoritmo aplica un \textbf{filtrado topol\'ogico a priori}: si dos celdas de la jaula residen en la misma fila o columna de la cuadr\'{\i}cula, cualquier tupla que repita un d\'{\i}gito entre ellas es podada antes de compilar el modelo. Esta formulaci\'on extensional permite al solver CP-SAT mantener \textbf{Consistencia de Arco Generalizada (GAC)}~\cite{bessiere2006table} sobre la totalidad del \textit{scope} de la jaula sin recurrir a variables auxiliares intermedias de multiplicaci\'on.

\subsection{Variante C: Inferencia Conjunta Neuro-Simb\'olica v\'{\i}a Reificaci\'on con Regularizaci\'on \texorpdfstring{$\ell_0$}{L0}}
\label{sec:joint_inference}
Para mitigar la propagaci\'on de errores sensoriales de la etapa de visi\'on, la Variante C transforma el CSP en un problema de optimizaci\'on neuro-simb\'olica de M\'axima Verosimilitud Consistente (MAP)~\cite{deraedt2011constraint}.

Para cada jaula $c$, la etapa de percepci\'on suministra un conjunto de $K_c$ hip\'otesis de lectura con sus correspondientes probabilidades normalizadas $P_{c,k} \in (0, 1]$. Se introducen variables booleanas de selecci\'on:
\begin{equation}
    r_{c,k} \in \{0, 1\}, \quad \forall c \in \mathcal{C}, \; \forall k \in \{0, \dots, K_c - 1\}
\end{equation}
sujetas a la restricci\'on de partici\'on de hip\'otesis \'unica:
\begin{equation}
    \text{ExactlyOne}(\{r_{c,0}, r_{c,1}, \dots, r_{c,K_c - 1}\}), \quad \forall c \in \mathcal{C}
\end{equation}

Las restricciones aritm\'eticas de cada candidata se vinculan condicionalmente mediante reificaci\'on v\'{\i}a \texttt{OnlyEnforceIf}:
\begin{equation}
    r_{c,k} \implies \mathcal{A}(V_c, T_{c,k}, \text{op}_{c,k})
\end{equation}

\subsubsection{Funci\'on Objetivo con Regularizaci\'on \texorpdfstring{$\ell_0$}{L0}}
En formulaciones iniciales sin regularizaci\'on, cuando varias hip\'otesis alternativas pose\'{\i}an probabilidades bajas similares, el solver permutaba m\'ultiples jaulas arbitrariamente para satisfacer el Cuadrado Latino, originando soluciones v\'alidas en l\'ogica pero divergentes de la imagen. Para erradicar este fen\'omeno, se incorpor\'o una \textbf{penalizaci\'on $\ell_0$ constante} $\lambda_{\text{change}} = 1.5$ ($1500$ unidades enteras normalizadas) que penaliza estrictamente cualquier desviaci\'on de la hip\'otesis visual Top-1 ($k > 0$):
\begin{equation}
\begin{aligned}
    \max \sum_{c \in \mathcal{C}} \sum_{k=0}^{K_c - 1} r_{c,k} \cdot \Big( & \lfloor 1000 \cdot \log P_{c,k} \rfloor \\
    & - \mathbf{1}_{\{k > 0\}} \cdot 1500 \Big)
\end{aligned}
\label{eq:map_objective_l0}
\end{equation}

El solver audita deterministamente la cantidad total de jaulas alteradas:
\begin{equation}
    M_{\text{corregidas}} = \sum_{c \in \mathcal{C}} \sum_{k=1}^{K_c - 1} r_{c,k}
\end{equation}
asegurando que las correcciones respondan a rescates leg\'{\i}timos de ambig\"uedades puntuales y no a mutaciones arbitrarias del acertijo original.


\section{Conclusiones y Trabajo Futuro}
La integraci\'on de t\'ecnicas de Programaci\'on por Restricciones (CP-SAT) con visi\'on por computador ha permitido desarrollar un sistema robusto de alta eficiencia para la resoluci\'on de KenKen. El uso de restricciones globales extensionales de tabla (\texttt{AllowedAssignments}) demostr\'o alcanzar Consistencia de Arco Generalizada (GAC), logrando una aceleraci\'on de casi el 40\% en $n=9$ respecto al modelo intensional. Asimismo, el esquema de inferencia conjunta mediante reificaci\'on probabil\'{\i}stica confiere al sistema una tolerancia completa frente a errores sensoriales del OCR.

Como trabajo futuro, se proyecta la extensi\'on del marco de inferencia hacia aprendizaje end-to-end de pesos de confianza mediante diferenciaci\'on de solvers combinatorios y la evaluaci\'on del pipeline sobre datasets reales capturados en condiciones lum\'{\i}nicas adversas.

\bibliographystyle{IEEEtran}
\bibliography{refs}

\end{document}
